%======================================================================
\subsection{Part properties that can grow (built 19 Oct 2026)}\label{sec:partprops}
%======================================================================

\textbf{Requirement.} ``Make it so we can expand properties of parts later.'' The properties every part carries
(P/N, revision, description, material, mass, finish, vendor\ldots) are therefore data, not code: a list in
JSON that is read in layers. Adding a property means adding one line to a layer file; FreeCAD, the BOM and the
vault pick it up without a new build.

\begin{center}
\small
\begin{tabularx}{\linewidth}{@{}lX@{}}
\toprule
Layer (read in this order) & Holds\\
\midrule
1. Built in: \file{addons/BruceWorx/config/part\_properties.json} & The ten Heuretes FreeCAD properties below. Ships with
  the build; not edited by hand.\\
2. Company or vault: e.g.\ \file{<vault>/part\_properties.json} & What your company adds: supplier, lead time,
  RoHS, cost centre\ldots\\
3. User: \file{\textasciitilde/.config/bruceworx/part\_properties.json} & Personal additions while trying something
  out.\\
\bottomrule
\end{tabularx}
\end{center}

\begin{center}
\small
\begin{tabular}{@{}lllll@{}}
\toprule
Property & Type & BOM column & Comes from & Needed at\\
\midrule
\code{PartNumber} & text & \code{.PartNumber} & vault (read-only) & ---\\
\code{Revision} & text & \code{.Revision} & vault (read-only) & ---\\
\code{LifecycleState} & choice & on request & vault (read-only) & ---\\
\code{Description} & text & \code{.Description} & you & submit\\
\code{ItemType} & choice & --- & you (part, assy, COTS, dwg) & ---\\
\code{MaterialName} & text & \code{.MaterialName} & computed: bodies' material & submit\\
\code{Finish} & text & \code{.Finish} & you & ---\\
\code{Mass} & mass & \code{.Mass} & volume $\times$ density & ---\\
\code{Vendor}, \code{VendorPart\-Number} & text & --- & you & ---\\
\bottomrule
\end{tabular}
\end{center}

\paragraph{Rules that keep it safe to grow.}
\begin{itemize}[nosep]
  \item A later layer may \emph{add} properties and may change an existing one's description, default, BOM flag
        or the stage that needs it. It may \emph{not} change a property's type; that would make old files
        unreadable, so the layer is refused with a message naming both files.
  \item Renames are declared (\code{"renamed": \{"Old": "New"\}}); applying the schema moves the value.
  \item Applying never deletes a property it does not know, so a file made with a bigger schema keeps its extra
        values when a smaller one opens it.
  \item A clash with a property FreeCAD already has is reported and left alone. The first real run found one:
        \code{App::Part} already has a \code{Material} link, so ours is \code{MaterialName}.
  \item ``Needed at'' feeds the lifecycle: \code{submit} lists what is still empty before a part can go to review.
\end{itemize}

\paragraph{Lifecycle state (built 22 Oct 2026).} A part's state is a property the vault sets, not a guess from
its revision code. \code{LifecycleState} is a choice of exactly the vault's states (\code{WIP}, \code{In Review},
\code{Released}, \code{Obsolete}; section~\ref{sec:pdm}), starts at \code{WIP} and is read-only in the property
editor. A layer may restyle it (description, BOM flag) but not change its states. Two paths write it, both through
one function, \code{partprops.set\_state}: a vault opened in FreeCAD tells every open part whose number changes
state (submit, reject, approve, an ECO opening or releasing, retire), and a document opened while
\code{BruceWorx/VaultPath} names a vault takes each known part's state from the vault's journal (also from
the command line: \code{bw-hooks vault sync -{}-file part.FCStd}). The vault model itself stays free of FreeCAD: it
only calls listeners with \code{\{P/N: state\}}. The property is not named \code{State}: that is FreeCAD's own
\code{obj.State} (Touched, Invalid\ldots), and a property of that name would hide it from every macro and
workbench that reads it; applying the schema now refuses any property that would hide an attribute. The Browser
badge reads the property (WIP in the warning colour, In Review in the accent, Released quiet, Obsolete faded) and
guesses from the revision only for parts without it. Switch: Preferences \textrightarrow{} BruceWorx \textrightarrow{}
\emph{Lifecycle state of parts follows the vault}.

\paragraph{BOM.} The Assembly BOM prints any property whose column is named \code{.Property}, so the schema
just lists those names. A choice prints its text, so \code{.LifecycleState} works as a column; it is not one of the
default columns. Running it for real exposed an upstream bug: the BOM spreadsheet turned the text
``01'' into the number 1, so revision 01 printed as ``1'' (and a P/N with leading zeros would lose them).
\file{BomObject.cpp} now stores text values as text, the same way it already stores the Index column.

\paragraph{Verified headless.} 12 unit tests (layers, renames, clashes, required stages, enum choices), and a
FreeCAD run on an aluminium plate: mass 135\,g from its 50\,000\,mm$^3$ and Aluminum-6061-T6 density; material
name taken from the body; the BOM row reads \code{LE-01-01001-001 | 01 | Base plate | Aluminum-6061-T6 | Clear
anodise | 135 g}. The Assembly test suite still passes.

\paragraph{Not yet.} A FreeCAD command and panel to apply and edit the properties (today: one line in the Python
console, step 17 of the checklist), and recomputing Mass whenever the body changes (today: whenever the schema is
applied). Both belong to \eref{ASM-04} and \eref{PDM-04}.
